\documentclass[12pt,a4paper]{article}

\usepackage[utf8]{inputenc}
\usepackage[T1]{fontenc}
\usepackage{lmodern}
\usepackage[margin=1in]{geometry}
\usepackage{setspace}
\usepackage[hidelinks]{hyperref}
\usepackage{amsmath,amssymb}
\usepackage{booktabs}
\usepackage{titlesec}
\usepackage[stretch=15]{microtype}
\usepackage{parskip}
\usepackage{longtable}
\usepackage[
	backend=biber,
	style=authoryear,
	maxcitenames=3,
	mincitenames=1,
	uniquelist=false,
	sorting=nyt,
	dashed=false,
]{biblatex}
\addbibresource{references.bib}

% Section formatting
\titleformat{\section}{\large\bfseries}{}{0em}{}
\titleformat{\subsection}{\normalsize\bfseries}{}{0em}{}
\titleformat{\subsubsection}{\normalsize\itshape}{}{0em}{}

\title{Technical Note: The RDU--CRRA--Prelec Model as a Workhorse\\[6pt]
\large A Short Literature Review Supporting the Claim in Section~4}
\author{}
\date{March 2026}

\begin{document}
\maketitle
\thispagestyle{empty}

\section{Purpose}

In Section~4 of the paper, after specifying the 4-parameter RDU model---CRRA
utility $u(x) = x^{1-r}/(1-r)$, the 2-parameter \textcite{Prelec1998}
probability weighting function $w(p) = \exp(-\eta(-\ln p)^\phi)$, and a Fechner
noise parameter $\lambda$ with the \textcite{Wilcox2011} contextual utility
normalization---the text states:

\begin{quote}
``This 4 parameter model is a common workhorse of structural models of risk
preferences, used in a varied of applications.''
\end{quote}

This note compiles the evidence supporting that claim, organized by application
domain. The model specification---RDU with CRRA utility, 2-parameter Prelec
PWF, and a stochastic error specification (typically Fechner or contextual
utility)---constitutes 4 individual-level parameters $(r, \phi, \eta, \lambda)$.
Variants that omit the noise parameter or use a 1-parameter PWF are noted where
they appear but are not the primary focus.

\section{The Model's Core Components}

Rank-Dependent Utility \parencite{Quiggin1982} generalizes Expected Utility
Theory by allowing decision weights to differ from objective probabilities
through a probability weighting function (PWF). The model nests EUT as a
special case when the PWF is linear ($\phi = \eta = 1$ in the Prelec
specification).

The combination of CRRA utility with the 2-parameter Prelec PWF has become the
dominant specification in structural estimation of risk preferences for several
reasons. First, CRRA is the most widely used parametric utility function in the
risk preference literature, owing to its single-parameter parsimony and the
scale invariance it implies \parencite{HarrisonRutstrom2008}. Second, the
\textcite{Prelec1998} 2-parameter function
$w(p) = \exp(-\eta(-\ln p)^\phi)$ is the only axiomatically derived PWF, and
it nests both the 1-parameter power function ($\phi = 1$) and linear
probability weighting ($\phi = \eta = 1$) as special cases. Empirical
comparisons across functional forms have found that the 2-parameter Prelec form
is preferred to alternatives such as the \textcite{TverskyKahneman1992}
specification \parencite{Stott2006, HarrisonSwarthout2023}. Third, the
\textcite{Wilcox2011} contextual utility model provides a theoretically
grounded heteroscedastic error specification that ensures estimated utility
parameters correspond to a ``stochastically more risk averse'' ordering, which
is not guaranteed by homoscedastic Fechner or Luce error models.

\section{Applications by Domain}

\subsection{Laboratory Estimation of Risk Preferences}

The foundational application of this model is the structural estimation of
individual risk preferences from incentivized lottery choice experiments.
\textcite{Harrison2007} conducted a large-scale field experiment in Denmark
with a representative sample of 253 subjects, estimating CRRA coefficients
under EUT and then extending to RDU with probability weighting. Their estimates
indicated an average CRRA coefficient of approximately 0.67 and could not
reject CRRA as a characterization of behavior across demographic subgroups.
This study established the methodological template that much subsequent work
follows: elicit binary lottery choices, estimate the structural parameters of a
stochastic choice model via maximum likelihood.

\textcite{Andersen2008}, in an influential paper published in
\textit{Econometrica}, designed experiments to jointly elicit risk and time
preferences for the adult Danish population. The risk preference component used
CRRA utility with probability weighting under RDU. Their central finding was
that accounting for the curvature of the utility function (i.e., risk aversion)
when estimating discount rates leads to substantially lower---and more
plausible---estimated discount rates than when utility curvature is ignored.
This paper demonstrated that the structural RDU model with CRRA utility is
necessary not just for studying risk preferences in isolation, but for any
domain where risk preferences confound other preference parameters.

\subsection{Bayesian and Hierarchical Estimation}

\textcite{Nilsson2011} implemented hierarchical Bayesian parameter estimation
for Cumulative Prospect Theory (CPT), using the CRRA--Prelec specification as
the core parametric model. Their approach estimates individual-level parameters
drawn from group-level distributions, yielding more reliable estimates than
individual maximum likelihood, especially for subjects with limited choice
data. This paper established the feasibility and advantages of Bayesian
hierarchical estimation for prospect-theoretic models and has been widely cited.

\textcite{Gao2022} proposed a full Bayesian hierarchical model for RDU with
CRRA utility and the 2-parameter Prelec PWF, applied to the normative
evaluation of insurance purchase decisions. Their model estimates
individual-level parameters $(r, \phi, \eta, \lambda)$ within a hierarchical
structure where individual parameters are drawn from population distributions.
They demonstrated that Bayesian estimation provides more systematic control
over informative priors for individual risk preferences, and that these priors
make a substantive difference to welfare evaluations. This is the direct
predecessor of the model specification used in the present paper.

\textcite{Harrison2025} extends the Bayesian hierarchical framework to
accommodate \textit{mixtures} of risk preference models---allowing individuals
to use some combination of EUT, RDU, and other models---while maintaining the
CRRA--Prelec parameterization as the core specification within the RDU
component.

\subsection{Behavioral Welfare Economics}

\textcite{Harrison2016} applied the RDU--CRRA--Prelec model to evaluate the
expected welfare gain from insurance products. They estimated individual risk
preferences from one experimental task and then used those estimated
preferences to compute the certainty equivalent of risky insurance policies,
comparing it to the certain insurance premium. Their finding---that the metric
of take-up itself, widely used in policy evaluation, provides a qualitatively
incorrect guide to welfare effects---depends critically on having a structural
model rich enough to distinguish probability weighting from utility curvature.

\textcite{HarrisonSwarthout2023} conducted a systematic comparison of EUT,
RDU, Dual Theory, and CPT using laboratory data with real gains and losses. The
CRRA--Prelec parameterization was used for all non-EUT models. Their central
finding was that RDU is the best single characterization of individual and
pooled choices, outperforming CPT, and that much of what has been attributed to
``loss aversion'' in the CPT literature is better explained by probability
weighting under RDU.

\subsection{Joint Estimation of Risk and Time Preferences}

\textcite{Andersen2014} used the RDU--CRRA--Prelec model to jointly estimate
risk and time preferences of small business entrepreneurs in a Danish field
experiment. The risk preference specification included the full 4-parameter
model, with CRRA utility and the 2-parameter Prelec PWF estimated under RDU.
They found that while entrepreneurs are not significantly more or less risk
averse than the general population under EUT, they do exhibit different
probability weighting patterns under RDU---specifically, greater optimism about
the best outcomes. This application illustrates the model's use in
characterizing heterogeneity across economically meaningful subpopulations.

\subsection{Subjective Belief Elicitation}

\textcite{Harrison2017} and \textcite{Harrison2022a} addressed the problem of
recovering subjective probability distributions from incentivized scoring rule
data. When an individual's preferences conform to RDU rather than EUT, their
reported beliefs under a proper scoring rule are systematically distorted by
probability weighting. The correction procedure requires structural estimation
of the RDU--CRRA--Prelec model from a separate risk preference task, and then
using those estimates to ``undo'' the distortion in reported beliefs. This
application demonstrates that the model is not merely a tool for studying risk
preferences in isolation, but is a necessary input to any elicitation procedure
where risk preferences confound observed reports.

\subsection{COVID-19 Pandemic and Background Risk}

\textcite{Harrison2022b} applied the RDU--CRRA--Prelec model to an online,
incentivized experiment conducted during the COVID-19 pandemic in the United
States. They estimated EUT and RDU models using the standard CRRA utility
function and 2-parameter Prelec PWF, finding that the atemporal risk premium
changed significantly during the pandemic compared to pre-pandemic
baselines---consistent with theoretical predictions about the effect of
increased background risk on foreground risk attitudes. This application
demonstrates the model's use in tracking preference stability (or instability)
over time in response to naturally occurring shocks.

\subsection{Field Experiments in Developing Countries}

Research on risk preferences in developing countries has also adopted the
RDU--CRRA--Prelec specification. \textcite{Tanaka2010} linked experimental and
household survey data from Vietnam, using RDU with CRRA utility and a
probability weighting function to jointly estimate utility curvature and
probability weighting from incentivized lottery experiments. They found that a
substantial fraction of subjects exhibit behavior consistent with RDU rather
than EUT, and that restricting estimation to EUT with CRRA poorly predicts
preferences for less risk-averse subjects. See also
\textcite{HarrisonRutstrom2008} for discussion of related field experiments in
developing countries using the same specification.

\subsection{Insurance Markets and Field Data}

\textcite{Barseghyan2013} used administrative data on insurance deductible
choices to estimate structural models of risky choice that incorporate both
``standard'' risk aversion (diminishing marginal utility, parameterized by
CRRA) and probability distortions (parameterized by a PWF). They found that
probability distortions---characterized by substantial overweighting of small
probabilities---play an important role in explaining insurance deductible
choices, beyond what utility curvature alone can account for.
\textcite{Barseghyan2018} provide a comprehensive survey of field estimation of
risk preferences using these models. While their specification uses a different
PWF parameterization than the Prelec form, the underlying model
structure---separating utility curvature from probability weighting within an
RDU framework---is the same.

Harrison and collaborators at CEAR have applied the full CRRA--Prelec--RDU
specification to evaluate insurance experiments, examining both the expected
welfare gain from insurance \parencite{Harrison2016} and the welfare effects of
insurance contract non-performance. \textcite{Harrison2019} documents the
extensive use of the RDU--CRRA--Prelec model in behavioral insurance economics.

\subsection{Mixture Models and Model Comparison}

\textcite{Conte2011}, published in the \textit{Journal of Econometrics}, used
mixture models to address heterogeneity in preference functionals across
individuals, comparing EUT and RDU specifications with CRRA utility. Their
approach allows individuals to probabilistically use different decision
theories across different choice problems, with the CRRA--PWF specification
serving as the parametric core.

\section{Summary Table of Applications}

\begin{longtable}{p{4.5cm} p{5cm} p{4.5cm}}
\toprule
\textbf{Domain} & \textbf{Key References} & \textbf{Application} \\
\midrule
\endfirsthead
\toprule
\textbf{Domain} & \textbf{Key References} & \textbf{Application} \\
\midrule
\endhead
Laboratory risk elicitation & \textcite{Harrison2007} & Representative sample field experiment \\[4pt]
Joint risk and time & \textcite{Andersen2008} & Correcting discount rate estimates \\[4pt]
Hierarchical Bayesian & \textcite{Nilsson2011}; \textcite{Gao2022} & Individual-level estimation with shrinkage \\[4pt]
Welfare economics & \textcite{Harrison2016} & Insurance welfare evaluation \\[4pt]
Model comparison & \textcite{HarrisonSwarthout2023} & EUT vs.\ RDU vs.\ CPT \\[4pt]
Entrepreneurship & \textcite{Andersen2014} & Entrepreneur risk/time preferences \\[4pt]
Belief elicitation & \textcite{Harrison2017}; \textcite{Harrison2022a} & Scoring rule correction for PWF \\[4pt]
Pandemic shocks & \textcite{Harrison2022b} & COVID-19 background risk effects \\[4pt]
Developing countries & \textcite{Tanaka2010} & Smallholder farmer risk attitudes \\[4pt]
Insurance field data & \textcite{Barseghyan2013}; \textcite{Barseghyan2018} & Deductible choice, PWF identification \\[4pt]
Mixture models & \textcite{Conte2011} & Heterogeneous preference functionals \\[4pt]
Mixture + hierarchical & \textcite{Harrison2025} & Bayesian mixture of risk models \\
\bottomrule
\end{longtable}

\section{Assessment}

The claim that the 4-parameter RDU--CRRA--Prelec model is a ``common workhorse
of structural models of risk preferences'' is well supported by the literature.
The model has been applied across at least a dozen distinct application domains
over nearly two decades, by research groups spanning at least three continents,
using both laboratory and field data, with both frequentist and Bayesian
estimation methods. The CRRA utility function is the most common parametric
utility specification in the literature, the 2-parameter Prelec PWF is the most
common and best-justified probability weighting function, and RDU is
increasingly recognized as the best single-model characterization of individual
risky choice \parencite{HarrisonSwarthout2023}.

The 4-parameter version---including the noise parameter $\lambda$---is
specifically prominent in the structural estimation tradition associated with
Harrison and collaborators at CEAR, which has produced the majority of applied
work using this exact specification. Other research groups
\parencite{Nilsson2011, Conte2011, Barseghyan2013} use closely related
specifications that share the same parametric core. The model's reach into
welfare economics, belief elicitation, joint preference estimation, and field
applications in developing countries further attests to its status as a
workhorse.

\section{Suggested Revision for the Paper}

The current sentence reads:

\begin{quote}
``This 4 parameter model is a common workhorse of structural models of risk
preferences, used in a varied of applications.''
\end{quote}

A possible revision that is more precise and adds citations:

\begin{quote}
``This 4-parameter model is a workhorse of structural estimation of risk
preferences, applied to welfare evaluation of insurance decisions (Harrison and
Ng, 2016), joint estimation of risk and time preferences (Andersen et al.,
2008), recovery of subjective belief distributions (Harrison et al., 2017),
hierarchical Bayesian estimation (Gao et al., 2022; Nilsson et al., 2011), and
field experiments in developing countries (Tanaka et al., 2010), among many
others.''
\end{quote}

This version replaces the vague ``used in a varied of applications'' (which
also contains a typo) with specific citations spanning distinct domains,
allowing the reader to verify the breadth claim. The citations are chosen to
span different research groups, journals, and application areas.

\printbibliography

\end{document}
